\section{Arithmetic pair information and its localization limits}\label{sec:arithmetic}

The full smoothed pair function retains a common center, a separation,
and two transverse marks. We now compare that information with one known
arithmetic projection. The comparison explains which variables an existing
theorem integrates out, and therefore which new estimates the present
program needs. It also supports the marked scalar specialization in
Appendix~\ref{sec:pairs}. The conclusions concern the stated imports;
they are not a limitation theorem for all arithmetic correlation methods.

\subsection{The corrected cumulative and dyadic imports}

Write $L=\log T$, $B_T=TL/(2\pi)$, and use $r_0>0$ for the
exponential base in the arithmetic formula. For an ordinate interval
$I$ define
\begin{equation}\label{eq:arithmetic-F}
 \mathcal F_I(r_0)=
 \sum_{\gamma_j,\gamma_k\in I}
       r_0^{\rho_j-\rho_k}W(\rho_j-\rho_k),
 \qquad W(u)=\frac4{4-u^2}.
\end{equation}
All zeta zeros and all multiplicity copies are included, along
with their self labels. The corrected source is
Baluyot--Goldston--Suriajaya--Turnage-Butterbaugh
\cite[Section 3]{BGST2026PairI}, version 3.
The first line of its cumulative formula (3.5) gives, uniformly
for $1\le r_0\le T$,
\begin{equation}\label{eq:bgst-cumulative-raw}
 \mathcal F_{(0,T]}(r_0)
 =\frac{T}{2\pi r_0^2}\{L^2+O(L)\}
       +\frac{T}{2\pi}\log r_0+O(T\sqrt L).
\end{equation}
Consequently, for $0\le\alpha\le1$,
\begin{equation}\label{eq:bgst-cumulative}
 \frac{\mathcal F_{(0,T]}(T^\alpha)}{B_T}
 =T^{-2\alpha}\{L+O(1)\}+\alpha+O(L^{-1/2}).
\end{equation}
The corresponding dyadic theorem, its equation (3.2), is
\begin{equation}\label{eq:bgst-dyadic}
 \frac{\mathcal F_{(T,2T]}(T^\alpha)}{B_T}
 =T^{-2\alpha}L+\alpha
       +O(T^{-2\alpha}\sqrt L+L^{-1/2}).
\end{equation}
The stronger first error in \eqref{eq:bgst-cumulative} must not
be transferred to the different population in
\eqref{eq:bgst-dyadic}. These are imported arithmetic estimates,
not reproved prime mean-square theorems in this draft. Their
uniform big-$O$ constants and onsets are not numerically
instantiated here.

For clarity, positivity of the full sum has an elementary finite
explanation. If $u_j=\rho_j-1/2$, the zero multiset in $I$ is
invariant under $u\mapsto-\overline u$. The identity
$W(u)=\int_\R\ee^{-2|t|}\ee^{tu}\dd t$ for $|\Re u|<2$
therefore gives
\begin{equation}\label{eq:arithmetic-gram}
 \mathcal F_I(r_0)=
 \int_\R\ee^{-2|t|}
       \left|\sum_{\gamma_j\in I}
                   \ee^{(\log r_0+t)u_j}\right|^2\dd t\ge0.
\end{equation}
Interchange is valid for the finite population and
$|\Re(\rho_j-\rho_k)|\le1$.
Swapping labels shows
$\mathcal F_I(r_0^{-1})=\mathcal F_I(r_0)$.
Neither observation makes every complex summand nonnegative.

\subsection{A triangular-frequency projection and its actual value}

Let $q_{\rm freq}(\alpha)=(1-|\alpha|)_+$ and define its entire
transform
\begin{equation}\label{eq:fejer-transform}
 K_T(u)=\int_{-1}^1(1-|\alpha|)\ee^{\alpha Lu}\dd\alpha
       =\left(\frac{\sinh(Lu/2)}{Lu/2}\right)^2,\qquad K_T(0)=1.
\end{equation}
One can verify the formula by integrating twice, or by observing
that the triangular function is the convolution of two unit
interval indicators and multiplying their exponential transforms.
The ordered factorial statistic for the dyadic population
$\mathcal Z(T)=\{\rho:T<\gamma\le2T\}$ is
\begin{equation}\label{eq:A_T}
 A_T=\frac1{B_T}\sum_{\substack{\rho_j,\rho_k\in\mathcal Z(T)\\j\ne k}}
       W(\rho_j-\rho_k)K_T(\rho_j-\rho_k).
\end{equation}
Writing $n_T=|\mathcal Z(T)|$, the removed self mass is exactly
$n_T/B_T$. Distinct labels at a multiple zero remain.

\begin{proposition}[Value of the imported projection]\label{prop:bulk-value}
The dyadic import \eqref{eq:bgst-dyadic} and the classical
unconditional count $n_T=B_T+O(T)$ imply
\begin{equation}\label{eq:bulk-value}
 A_T=\frac13+O(L^{-1/2}).
\end{equation}
If the error in \eqref{eq:bgst-dyadic} is bounded by
$C_{\rm BG}(T^{-2\alpha}\sqrt L+L^{-1/2})$ and
$|n_T/B_T-1|\le C_N/L$ in a common range, then
\begin{equation}\label{eq:bulk-error}
 \left|A_T-\frac13\right|
 \le\frac{2C_{\rm BG}}{\sqrt L}+\frac{C_N+1/2}{L}.
\end{equation}
\end{proposition}
\begin{proof}
Integrate \eqref{eq:bgst-dyadic} against the triangular weight,
use evenness, and subtract the self mass. The two leading integrals
are
\[
 2L\int_0^1(1-\alpha)\ee^{-2L\alpha}\dd\alpha
       =1-\frac{1-T^{-2}}{2L},\qquad
 2\int_0^1(1-\alpha)\alpha\dd\alpha=\frac13.
\]
The integrated first error is at most
$2C_{\rm BG}\sqrt L\int_0^1\ee^{-2L\alpha}\dd\alpha
\le C_{\rm BG}/\sqrt L$.
The second contributes $C_{\rm BG}/\sqrt L$.
The zero count and the remaining $1/(2L)$ term give
\eqref{eq:bulk-error}. In particular the frequency-zero spike
has been retained.
\end{proof}

This is an actual arithmetic value for the declared projection.
It is not the value of the marked test in Appendix~\ref{sec:pairs}.
In the $H$ coordinate, the population lies in $2T<x\le4T$,
a window of half-width $T$ centered at $3T$.
Moreover
\[
 L(\rho_j-\rho_k)=
 \frac L2\bigl(-(y_j-y_k)+\ii(x_j-x_k)\bigr),
\]
so the microscopic difference scale is $4\pi/L$ in the $x$
coordinate. The projection retains a complex difference but
neither arbitrary common-center localization nor the two
separate transverse marks.

A simple $H$-pair at $x_*\pm\ii d$, $0<d<1$, has zeta
real-part difference $d$. Its two ordered mutual terms in
\eqref{eq:A_T} contribute
\begin{equation}\label{eq:arithmetic-defect}
 \frac2{B_T}\frac4{4-d^2}
       \left(\frac{\sinh(dL/2)}{dL/2}\right)^2 .
\end{equation}
For $d=c/L$ with fixed $c>0$, this is $O((TL)^{-1})$,
while the available error allowance in \eqref{eq:bulk-error}
is of order $L^{-1/2}$ and the main value is $1/3$.
At $d=0$ two different coincident real labels contribute $2/B_T$.
Other complex summands need not be nonnegative, so
\eqref{eq:arithmetic-defect} is neither a lower bound on $A_T$
nor the total change produced by inserting a pair: cross
interactions remain.
Thus this known projection supplies neither the desired
short-window marked value nor a useful absolute single-defect
comparison with the stated error.

\subsection{What the real-root and replacement-method benchmarks supply}

The classical Montgomery pair-correlation framework
\cite{Montgomery1973} motivates distinguishing bulk difference
information from a fully marked local marginal. An elementary
real-root identity illustrates the distinction. For a polynomial
with real roots $a_1,\ldots,a_n$, counted with multiplicity, put
\[
 p_P(x)=\sum_{j=1}^n\frac{y}{\pi((x-a_j)^2+y^2)},\qquad y>0 .
\]
The Cauchy densities satisfy $P_y*P_y=P_{2y}$; this follows,
for example, from their Fourier transform $\ee^{-y|\xi|}$.
Therefore
\begin{equation}\label{eq:real-poisson}
 2\pi y\norm{p_P}_{L^2(\R)}^2-n
 =\sum_{j\ne k}\frac{4y^2}{(a_j-a_k)^2+4y^2}.
\end{equation}
Indeed the integral of the $j,k$ product is
$P_{2y}(a_j-a_k)$, and each self term becomes one after
multiplication by $2\pi y$.
At $y=h/2$ the separation kernel is exactly $c_h$.
This real-population benchmark, used in the earlier heat-flow
collision investigation, does not establish that an
approximation to the canonical field has only real roots.
Discarding its possible complex zeros would remove the very
information under investigation.

Likewise the local replacement principle of Tao and Vu
\cite[Theorems 2.1 and 3.1]{TaoVu2014LocalZeros}
is useful as an architecture, not an imported theorem about
$H_0$. Its hypotheses include nondegeneracy, non-clustering,
and joint smooth-test comparison of logarithmic values, in a
degree-at-most-$n$ polynomial setting; the real version also
requires reference-field repulsion bounds.
A compact-source entire transform is not a polynomial, and
its local Jensen count is not its polynomial degree.
Theorem~\ref{thm:pair-transfer} avoids that probabilistic
comparison problem for deterministic accuracy on one disk,
but supplies no comparison in law or favorable reference value.

The earlier Rodgers--Tao-related local heat-flow work, as
discussed in the predecessor comparison \cite{OCArithmetic},
supplies a further methodological lesson: real simple populations,
exterior forces, and same-root correlation errors must retain
their exact hypotheses and combined accounting. None of those
real-population hypotheses is produced here for a field with
possibly nonreal zeros. The separate compact-window argument in Section~\ref{sec:newman-time}
now gives a time modulus without a real-simple-population premise.
It does not supply a closed signed correlation hierarchy.
